\documentclass[11pt]{article}
\usepackage[margin=1in]{geometry}
\usepackage{amsmath,amssymb,amsthm}
\usepackage{microtype}
\DeclareMathOperator{\tr}{tr}

\title{Contrastive evidence, superposition, and our results:\\ a compatibility check}
\date{}

\begin{document}
\maketitle

\noindent\textbf{Setup.} Under our empirical finding $n=2$, the response has the form
\[
R(v) = F(v^\top A v), \qquad A = \sum_{i=1}^M w_i\, d_i d_i^\top, \quad w_i \geq 0,
\]
with feature directions $d_i \in \mathbb{R}^D$, $\|d_i\|=1$, and (under superposition) $M > D$. Sensitivity of a direction $u$ is $u^\top A u$.

\medskip
\noindent\textbf{Contrastive claim.} By isolating concepts via contrastive pairs and observing they steer the model more than random directions do, we have
\[
d_j^\top A d_j \;>\; \mu \quad \text{for every feature } j, \qquad \text{where } \mu = \tfrac{1}{D}\tr(A)
\]
is the average sensitivity of a uniformly random unit direction (i.e.\ $\mathbb{E}_r [r^\top A r] = \tr(A)/D$).

\medskip
\noindent\textbf{Question.} Is this compatible with our results and with superposition ($M > D$)?

\medskip
\noindent\textbf{Summing the inequality.} Multiply both sides by $w_j \geq 0$ (preserves the direction of the inequality) and sum over $j$:
\[
\sum_{j=1}^M w_j\, d_j^\top A d_j \;>\; \mu \sum_{j=1}^M w_j.
\]

\smallskip
\noindent\emph{Left side.} Using $d_j^\top A d_j = \tr(A\, d_j d_j^\top)$ and linearity of trace,
\[
\sum_j w_j\, d_j^\top A d_j
= \tr\!\Big(A \sum_j w_j d_j d_j^\top\Big)
= \tr(A \cdot A)
= \tr(A^2).
\]

\smallskip
\noindent\emph{Right side.} Since $\|d_j\|=1$ implies $\tr(d_j d_j^\top) = 1$,
\[
\sum_j w_j = \sum_j w_j \tr(d_j d_j^\top) = \tr\!\Big(\sum_j w_j d_j d_j^\top\Big) = \tr(A).
\]

\smallskip
Substituting $\mu = \tr(A)/D$:
\[
\boxed{\tr(A^2) \;>\; \frac{\tr(A)^2}{D}.}
\]

\medskip
\noindent\textbf{What this says about $A$.} By Cauchy--Schwarz on the eigenvalues of $A$,
\[
\tr(A^2) = \sum_k \lambda_k^2 \;\geq\; \frac{(\sum_k \lambda_k)^2}{D} = \frac{\tr(A)^2}{D},
\]
with equality iff all eigenvalues are equal, i.e.\ $A \propto I$ (a tight frame). Our strict inequality just says \emph{$A$ is not a scaled identity} --- its eigenvalues are not all equal. This is generic and easily satisfied under $M > D$. \textbf{No contradiction.}

\medskip
\noindent\textbf{But: a non-feature direction must beat every feature.} Once $A$ has unequal eigenvalues, let $u^*$ be its top eigenvector. By Rayleigh,
\[
u^{*\top} A u^* = \lambda_{\max}(A) \;\geq\; d_j^\top A d_j \quad \text{for every } j.
\]
Under $M > D$ with generic features, $u^*$ does not coincide with any $d_j$ (that would be a measure-zero alignment among $M$ directions in $D$-space). So there exists a \emph{non-feature} direction strictly more sensitive than every feature.

\medskip
\noindent\textbf{Interpretation.} The contrastive evidence ``every feature beats a typical direction'' is fully consistent with our $n=2$ results and with superposition. What it \emph{does not} rule out is the existence of a non-feature direction that beats every feature --- in fact our results \emph{force} such a direction to exist (generically). A contradiction with the strong reading of LRH only arises if one separately asserts that no non-feature direction can exceed the sensitivity of every feature. The contrastive argument alone doesn't establish that: it compares features to typical directions, not to optimized ones.

\medskip
\noindent\textbf{Empirical closure.} The missing piece is direct: if one maximizes sensitivity over unit directions (MELBO-style) and the optimum consistently lands on recognizable feature directions, the strong reading holds. If the optimum lands on non-feature compositional directions, the strong reading fails --- which is the empirical form of the corollary, without any universal-quantifier assumption.

\end{document}
